\documentclass[12pt,a4paper,]{ctexart}
\usepackage[UTF8]{ctex}
\usepackage{amsmath,amssymb}
\usepackage{geometry}
\usepackage{xcolor}
\usepackage{listings}
\geometry{margin=2.5cm}
\title{Rose number}
\author{}
\date{}
\lstset{
	language=C++,
	basicstyle=\ttfamily\small,
	numbers=left,
	frame=single,
	tabsize=4,
	breaklines=true,
	showstringspaces=false,
	keywordstyle=\bfseries\color{red},
	directivestyle=\color{green},
	commentstyle=\color{green!70!black},
	stringstyle=\color{blue},
	numberstyle=\tiny\color{gray}
}
\begin{document}
	\maketitle
	\begin{center}
		Rose number(Four-leaf Rose Number)[四葉玫瑰数]\\
		Rose number is a Narcissistic number (自幂数) or Armstrong number (阿姆斯壯数)
	\end{center}
	This refers to a 4-digit number where the sum of the fourth powers of its digits equals the number itself.\\
	Example: 1634\\
	$1^4+6^4+3^4+4^4 = 1634$\\
	So \underline{1634} is a Four-leaf Rose Number.\\
	There have 3 rose number is:
	\begin{enumerate}
		\item 1634
		\item 8208
		\item 9474
	\end{enumerate}
	\newpage
	So it have a problem that is when we forget those number, can we find that?\\
	Yes, I some of the method to find those number.\\
	\section {Find Four-leaf Rose Number with programming.}
	\subsection{In C++}
	\begin{lstlisting}
		#include<iostream>
		#include<cstdio>
		using namespace std;
		
		int main(){
			int a,b,c,d;
			for (int num=1000;num<=9999;num++){
				a = num / 1000;
				b = (num / 100) % 10;
				c = (num / 10) % 10;
				d = num % 10;
				if (a*a*a*a+b*b*b*b+c*c*c*c+d*d*d*d==num)
				printf("%d\n",num);
			}
			cin.git();
			return 0;
		}
	\end{lstlisting}
	\subsection{In python}
	\begin{lstlisting}
		for num in range(1000,10000):
			a = num // 1000
			b = (num // 100) % 10
			c = (num // 10) % 10
			d = num % 10
			if a**4+b**4+c**4+d**4==num:
				print (num)
	\end{lstlisting}
\end{document}
